% Part: proof-theory
% Chapter: sequent-calculus
% Section: invertibility

\documentclass[../../../include/open-logic-section]{subfiles}

\begin{document}

\olfileid{pt}{seq}{inv}

\olsection{في قابلية القواعد للعكس}

\begin{defn}
نقول إن القاعدة~$R$،
\[
\AxiomC{$S_1$}
\AxiomC{$\dots$}
\AxiomC{$S_n$}
\RightLabel{$R$}
\TrinaryInfC{$S$}
\DisplayProof
\]
\emph{قابلة للعكس} في نظام ما متى كان ثبوت~$\Proves S$ يستلزم
$\Proves S_i$ لكل~$i=1$، \dots،~$n$. فإن بقي الارتفاع محفوظًا،
أي كان~$\Proves_n S$ مستلزمًا لـ$\Proves[n] S_i$ لكل~$i=1$، \dots،~$n$،
سميناها~\emph{قابلة للعكس مع حفظ~!!{height}}
(أو~\emph{قابلة للعكس مع حفظ~!!{hp}}).
\end{defn}

\begin{lem}\ollabel{lem:G3c-invert}
قواعد~\Log{G3c} المختصة بـ$\lnot$ و$\land$ و$\lor$ و$\lif$ قابلة للعكس
مع حفظ~!!{height}؛ وتفصيل ذلك ما يأتي:
\begin{enumerate}
  \item \RightR{\lnot}: إذا كان~$\Log{G3c} \Proves[n] \Gamma \Sequent
  \Delta, \lnot !A$، فإن~$\Log{G3c} \Proves[n] !A, \Gamma \Sequent
  \Delta$.
  \item \LeftR{\lnot}: إذا كان~$\Log{G3c} \Proves[n] \Gamma, \lnot !A
  \Sequent \Delta$، فإن~$\Log{G3c} \Proves[n] \Gamma \Sequent \Delta,
  !A$.
  \item \RightR{\land}: إذا كان~$\Log{G3c} \Proves[n] \Gamma \Sequent
  \Delta, !A \land !B$، فإن~$\Log{G3c} \Proves[n] \Gamma \Sequent
  \Delta, !A$ و$\Log{G3c} \Proves[n] \Gamma \Sequent
  \Delta, !B$.
  \item \LeftR{\land}: إذا كان~$\Log{G3c} \Proves[n] \Gamma, !A \land !B
  \Sequent \Delta$، فإن~$\Log{G3c} \Proves[n] !A, !B, \Gamma \Sequent
  \Delta$.
  \item \RightR{\lor}: إذا كان~$\Log{G3c} \Proves[n] \Gamma \Sequent
  \Delta, !A \lor !B$، فإن~$\Log{G3c} \Proves[n] \Gamma \Sequent
  \Delta, !A, !B$.
  \item \LeftR{\lor}: إذا كان~$\Log{G3c} \Proves[n] !A \lor !B, \Gamma \Sequent
  \Delta$، فإن~$\Log{G3c} \Proves[n] !A, \Gamma \Sequent
  \Delta$ و$\Log{G3c} \Proves[n] !B, \Gamma \Sequent
  \Delta$.
  \item \RightR{\lif}: إذا كان~$\Log{G3c} \Proves[n] \Gamma \Sequent
  \Delta, !A \lif !B$، فإن~$\Log{G3c} \Proves[n] !A, \Gamma \Sequent
  \Delta, !B$.
  \item \LeftR{\lif}: إذا كان~$\Log{G3c} \Proves[n] !A \lif !B, \Gamma
  \Sequent \Delta$، فإن~$\Log{G3c} \Proves[n] \Gamma \Sequent \Delta,
  !A$ و$\Log{G3c} \Proves[n] !B, \Gamma \Sequent \Delta$.
\end{enumerate}
\end{lem}

\begin{proof}
المطلوب في كل قاعدة~$R$ من~\Log{G3c} أن نستخرج، من كل~!!a{proof}~$\pi$
لنتيجتها~$S$ بحسب القاعدة~$R$، !!{proof}s~$\pi_i$ لمقدماتها~$S_i$، مع
$\pheight{\pi_i}\leq\pheight{\pi}$. ونفصّل الحجة في القاعدة
\LeftR{\land}. فليكن~$\pi$~!!a{proof} متتاليته النهائية
$!A\land !B,\Gamma\Sequent\Delta$؛ ونريد إنشاء~!!a{proof}~$\pi'$
لـ$!A,!B,\Gamma\Sequent\Delta$ يحقق~$\pheight{\pi'}\leq
\pheight{\pi}$. والسبيل إلى ذلك الاستقراء على~$n=\pheight{\pi}$.

أما إذا كان~$n=0$، فلا يعدو~$\pi$ بديهية واحدة أو تطبيقًا خالي المقدمات
للقاعدة~$\LeftR{\lfalse}$. وبديهيات~\Log{G3c} على الصورة~$!C,\Gamma
\Sequent\Delta,!C$، وشرطها أن تكون~$!C$ ذرية؛ فلا تكون إذن~$!C$ هي
$!A\land !B$. فإن كانت~$!A\land !B,\Gamma\Sequent\Delta$ بديهية، كان
بعض~$!C$ الذرية بوصفها~!!a{element} في كل من~$\Gamma$ و$\Delta$؛ ولذلك
تكون~$!A,!B,\Gamma\Sequent\Delta$ بديهية كذلك، وهي~!!{proof}~$\pi'$
ذو~!!{height}~$0$. وإن كان~$\pi$ تطبيقًا لـ$\LeftR{\lfalse}$، بقيت
$\lfalse$ في السياق الأيسر بعد الإحلال المذكور، فكانت
المتتالية المطلوبة نتيجة للقاعدة نفسها.

وأما إذا كان~$n>0$، فإن آخر البرهان استدلال. ونثبت الدعوى عند~$n$ بعد أن
نفترضها لكل~!!{proof} ارتفاعه دون~$n$، مهما اختلف سياقه. ومعنى الفرض
دقيقًا أنه لكل~$k<n$ ولكل~$\Pi$ و$\Lambda$، متى وُجد لـ$!A\land
!B,\Pi\Sequent\Lambda$~!!a{proof} ذو~!!{height}~$k$، وُجد~!!a{proof}
ذو~!!{height} لا يتجاوز~$k$ لـ$!A,!B,\Pi\Sequent\Lambda$.

والأمر على قسمين: إما أن تكون~$!A\land !B$ رئيسة في اليسار، وإما أن
تكون من السياق. فإن كانت رئيسة، كان الاستدلال الأخير~\LeftR{\land}،
وكانت خاتمة الـ!!{proof} الجزئي المباشر~$\pi_1$ هي
$!A,!B,\Gamma\Sequent\Delta$. فيكفينا ذلك البرهان نفسه، إذ
$\pheight{\pi_1}<\pheight{\pi}$.

وإن لم تكن~$!A\land !B$ رئيسة، كانت الرئيسة~!!{formula} أخرى، وكانت
$!A\land !B$ من سياق الاستدلال الأخير. فنطبق فرض الاستقراء على
الـ!!{proof}s الجزئية المباشرة للقاعدة الأخيرة~$R$؛ فينتج~!!a{proof}
واحد~$\pi_1'$، أو~!!{proof}s اثنان~$\pi_1'$ و$\pi_2'$، ولا يزيد
!!{height} واحد منها على ارتفاع نظيره من الـ!!{proof}s الجزئية المباشرة
لـ$\pi$. وخاتمة كل من~$\pi_1'$ و$\pi_2'$ تحمل~$!A,!B$ في السياق الأيسر
مكان~$!A\land !B$. ثم نعيد تطبيق~$R$ فنحصل على~$\pi'$. ولتوضيح ذلك،
لنفترض أن خاتمة~$\pi$ هي استدلال~\LeftR{\lif}:
\[
\AxiomC{}
\RightLabel{$\pi_1$}
\Deduce$!A \land !B, \Gamma' \fCenter \Delta, !C$
\AxiomC{}
\RightLabel{$\pi_2$}
\Deduce$!A \land !B, \Gamma', !D \fCenter \Delta$
\RightLabel{\LeftR{\lif}}
\BinaryInf$!A \land !B, \Gamma', !C \lif !D \fCenter \Delta$
\DisplayProof
\]
فالرئيسة في اليسار هنا~$!C\lif !D$، وسياقه الأيسر~$!A\land
!B,\Gamma'$؛ أي إن~$\Gamma=\Gamma',!C\lif !D$. ويفيد فرض الاستقراء
!!{proof}s~$\pi_1'$ و$\pi_2'$ لـ$!A,!B,\Gamma'\Sequent\Delta,!C$
و$!A,!B,\Gamma',!D\Sequent\Delta$ على الترتيب. ونجعل خاتمتيهما مقدمتين
جديدتين للقاعدة~\LeftR{\lif}، فيتم لنا~$\pi'$:
\[
\AxiomC{}
\RightLabel{$\pi_1'$}
\Deduce$!A, !B, \Gamma' \fCenter \Delta, !C$
\AxiomC{}
\RightLabel{$\pi_2'$}
\Deduce$!A, !B, \Gamma', !D \fCenter \Delta$
\RightLabel{\LeftR{\lif}}
\BinaryInf$!A, !B, \Gamma', !C \lif !D \fCenter \Delta$
\DisplayProof
\]
وحينئذ لدينا
\begin{align*}
\pheight{\pi} & = \max(\pheight{\pi_1}, \pheight{\pi_2}) + 1\\
\pheight{\pi'} & = \max(\pheight{\pi_1'}, \pheight{\pi_2'}) + 1
\intertext{بحسب التعريف، كما أن}
\pheight{\pi_1'} & \le \pheight{\pi_1}\\
\pheight{\pi_2'} & \le \pheight{\pi_2}
\intertext{بفرض الاستقراء. ومن ثم}
\pheight{\pi'} & \le \pheight{\pi}.
\end{align*}
\end{proof}

\begin{prob}
أعط برهان~\olref{lem:G3c-invert} للقاعدة~\RightR{\lif}.
\end{prob}

\begin{lem}\ollabel{lem:invert-quant}
القاعدتان~\RightR{\lforall} و\LeftR{\lexists} قابلتان للعكس مع
حفظ~!!{height} بالمعنى الآتي:
\begin{enumerate}
  \item إذا كان~$\Log{G3c} \Proves[n] \Gamma \Sequent \Delta,
  \lforall[x][!A(x)]$، فإن~$\Log{G3c} \Proves[n] \Gamma \Sequent
  \Delta,!A(c)$؛ و
  \item إذا كان~$\Log{G3c} \Proves[n] \lexists[x][!A(x)],\Gamma\Sequent
  \Delta$، فإن~$\Log{G3c} \Proves[n] !A(c),\Gamma\Sequent\Delta$،
\end{enumerate}
لكل~$c$ لا يظهر في~$\Gamma$ ولا~$\Delta$ ولا في~$\lforall[x][!A(x)]$
أو~$\lexists[x][!A(x)]$، على الترتيب.
\end{lem}

\begin{proof}
نقيم البرهان على (1)، وندع (2) تمرينًا؛ والحجة من جنس حجة
\olref{lem:G3c-invert}. ففي خطوة الاستقراء يعود التقسيم نفسه. فإن كانت
$\lforall[x][!A(x)]$ رئيسة، كانت مقدمة استدلال~$\RightR{\lforall}$
الأخير هي المتتالية المطلوبة~$\Gamma\Sequent\Delta,!A(a)$، وعندها ينتهي
الـ!!{proof} الجزئي المباشر~$\pi_1$، مع
$\pheight{\pi_1}<\pheight{\pi}$. ثم إن شرط المتغير المميز يضمن ألا يظهر
$a$ في~$\Gamma$ أو~$\Delta$ أو~$\lforall[x][!A(x)]$. ولنا أن نأخذ
$\pi_1$ منتظمًا؛ فيكون~$\Subst{\pi_1}{c}{a}$، بموجب
\olref[qua]{lem:subst-regular}، !!a{proof} من !!{height} نفسه
لـ$\Gamma\Sequent\Delta,!A(c)$.

وأما القسم الثاني فلا تكون فيه~$\lforall[x][!A(x)]$ رئيسة في الاستدلال
الأخير، بل تكون الرئيسة صيغة أخرى. ونكتفي بصورة ممثلة: ليكن الاستدلال
الأخير~\RightR{\land}. فتكون في~$\Delta$~!!{formula} من الصورة
$!B\land !C$ هي الرئيسة، وتكون خاتمة~!!{proof}~$\pi$ هكذا:
\[
\AxiomC{}
\RightLabel{$\pi_1$}
\Deduce$\Gamma \fCenter \Delta', \lforall[x][!A(x)], !B$
\AxiomC{}
\RightLabel{$\pi_2$}
\Deduce$\Gamma \fCenter \Delta', \lforall[x][!A(x)], !C$
\RightLabel{\RightR{\land}}
\BinaryInf$\Gamma \fCenter \Delta', \lforall[x][!A(x)], !B \land !C$
\DisplayProof
\]
حيث~$\Delta=\Delta',!B\land !C$. واختر~$c$~!!a{constant} لا تظهر في
$\Gamma$ ولا~$\Delta$ ولا~$\lforall[x][!A(x)]$؛ فلا تظهر، من باب أولى،
في~$\Delta',!B$ ولا في~$\Delta',!C$. فيجري فرض الاستقراء على~$\pi_1$
و$\pi_2$، ونحصل على~!!{proof}s~$\pi_1'$ و$\pi_2'$ مع
$\pheight{\pi_1'}\leq\pheight{\pi_1}$ و$\pheight{\pi_2'}\leq
\pheight{\pi_2}$. وأما الـ!!{proof} المطلوب~$\pi'$ لـ$\Gamma\Sequent
\Delta,!A(c)$ فهو:
\[
\AxiomC{}
\RightLabel{$\pi_1'$}
\Deduce$\Gamma \fCenter \Delta', !A(c), !B$
\AxiomC{}
\RightLabel{$\pi_2'$}
\Deduce$\Gamma \fCenter \Delta', !A(c), !C$
\RightLabel{\RightR{\land}}
\BinaryInf$\Gamma \fCenter \Delta', !A(c), !B \land !C$
\DisplayProof
\]
ومعه~$\pheight{\pi'}=\max(\pheight{\pi_1'},\pheight{\pi_2'})+1
\leq\pheight{\pi}$.
\end{proof}

لمّا كانت~\Olref{lem:G3c-invert} عمدةً في برهان مبرهنة حذف القطع، كان لها
أيضًا الأثر الآتي:

\begin{prop}\ollabel{prop:G3c-cont-adm}
قاعدتا التقلص~\LeftR{\Contraction} و\RightR{\Contraction} في~\Log{G1c}
مقبولتان في~\Log{G3c} مع حفظ~!!{height}.
\end{prop}

\begin{proof}
نسوق حجة واحدة لـ\RightR{\Contraction} و\LeftR{\Contraction}. فليكن
$\pi$~!!a{proof} في~\Log{G3c} لإحدى المتتاليتين
$\Gamma\Sequent\Delta,!A,!A$ أو~$!A,!A,\Gamma\Sequent\Delta$. والمطلوب
\Log{G3c}-!!{proof}~$\pi'$، على الترتيب، لـ$\Gamma\Sequent\Delta,!A$
أو~$!A,\Gamma\Sequent\Delta$، بحيث~$\pheight{\pi'}\leq
\pheight{\pi}$. ونستقرئ على~$n=\pheight{\pi}$.

فإن كان~$n=0$، كان~$\pi$ بديهية أو تطبيقًا بلا مقدمات لـ$\LeftR{\lfalse}$.
فإذا كانت~$\Gamma\Sequent\Delta,!A,!A$ بديهية لـ\Log{G3c}، كانت
$\Gamma\Sequent\Delta,!A$ بديهية أيضًا؛ إذ إما أن يوجد~$!C$ ذري
بوصفه~!!a{element} في كل من~$\Gamma$ و$\Delta$، وإما أن تكون~$!A$ ذرية
و!!a{element} في~$\Gamma$. وكذلك الحكم إذا كانت خاتمة البرهان
$!A,!A,\Gamma\Sequent\Delta$. وإن كان~$\pi$ تطبيقًا لـ$\LeftR{\lfalse}$،
ظلت~$\lfalse$ في الطرف الأيسر بعد التقلص، فنتج المطلوب بالقاعدة نفسها.

وإن كان~$n>0$، انتهى~$\pi$ بقاعدة~$R$. فإذا لم تكن~$!A$ رئيسة في هذا
الاستدلال، أنشأنا~!!a{proof}~$\pi'$ كما صنعنا في برهان
\olref{lem:G3c-invert}: نجري فرض الاستقراء على الـ!!{proof}s الجزئية
المباشرة لـ$\pi$، ثم نعيد القاعدة~$R$ على الـ!!{proof}s الناتجة. والفرض هنا عام؛
فمتى وُجد~!!a{proof} من~!!{height}~$k<n$ لإحدى المتتاليتين
$\Pi\Sequent\Lambda,!D,!D$ أو~$!D,!D,\Pi\Sequent\Lambda$، وُجد
!!a{proof} من~!!{height} لا يتجاوز~$k$ لـ$\Pi\Sequent\Lambda,!D$
أو~$!D,\Pi\Sequent\Lambda$ على الترتيب. فهو جار وإن خالفت السياقات أو
!!{formula} المتقلصة~$\Gamma$ و$\Delta$ و$!A$.

ومثال ذلك أن تكون~$R$ هي~$\RightR{\land}$، وأن ينتهي~$\pi$ بـ
\[
\AxiomC{}
\RightLabel{$\pi_1$}
\Deduce$\Gamma \fCenter \Delta', !B, !A, !A$
\AxiomC{}
\RightLabel{$\pi_2$}
\Deduce$\Gamma \fCenter \Delta', !C, !A, !A$
\RightLabel{\RightR{\land}}
\BinaryInf$\Gamma \fCenter \Delta', !B \land !C, !A, !A$
\DisplayProof
\]
حيث~$\Delta=\Delta',!B\land !C$. وعندئذ يفيد فرض الاستقراء~!!{proof}s
$\pi_1'$ و$\pi_2'$ لـ$\Gamma\fCenter\Delta',!B,!A$ و$\Gamma\fCenter
\Delta',!C,!A$ على الترتيب، مع~$\pheight{\pi_1'}\leq
\pheight{\pi_1}$ و$\pheight{\pi_2'}\leq\pheight{\pi_2}$. فيكون
الـ!!{proof}~$\pi'$:
\[
\AxiomC{}
\RightLabel{$\pi_1'$}
\Deduce$\Gamma \fCenter \Delta', !B, !A$
\AxiomC{}
\RightLabel{$\pi_2'$}
\Deduce$\Gamma \fCenter \Delta', !C, !A$
\RightLabel{\RightR{\land}}
\BinaryInf$\Gamma \fCenter \Delta', !B \land !C, !A$
\DisplayProof
\]

وأما إن كانت~$!A$ رئيسة في الاستدلال الأخير، بدأنا بلمّة العكس
(\olref{lem:G3c-invert}) على الـ!!{proof}s الجزئية المباشرة لـ$\pi$،
ثم أجرينا فرض الاستقراء، ثم أعدنا القاعدة~$R$. فافترض، مثلًا، أن~$\pi$
يبرهن~$\Gamma\Sequent\Delta,!A,!A$، وأن~$!A\ident(!B\land !C)$ هي
الرئيسة في استدلاله الأخير. حينئذ تكون صورة~$\pi$:
\[
\AxiomC{}
\RightLabel{$\pi_1$}
\Deduce$\Gamma \fCenter \Delta, !B \land !C, !B$
\AxiomC{}
\RightLabel{$\pi_2$}
\Deduce$\Gamma \fCenter \Delta, !B \land !C, !C$
\RightLabel{\RightR{\land}}
\BinaryInf$\Gamma \fCenter \Delta, !B \land !C, !B \land !C$
\DisplayProof
\]
فمن تطبيق~\olref{lem:G3c-invert} على~$\pi_1$ نحصل على~!!a{proof}~$\pi_1'$
لـ$\Gamma\Sequent\Delta,!B,!B$؛ وهي تعطينا أيضًا~!!a{proof} لـ$\Gamma
\Sequent\Delta,!C,!B$ لا حاجة لنا به. وبالمثل، يخرج من تطبيق
\olref{lem:G3c-invert} على~$\pi_2$~!!a{proof}~$\pi_2'$ لـ$\Gamma
\Sequent\Delta,!C,!C$. وإذ كان عكس~$\RightR{\land}$ حافظًا لـ!!{height}، فإن
$\pheight{\pi_1'}\leq\pheight{\pi_1}<\pheight{\pi}$ و
$\pheight{\pi_2'}\leq\pheight{\pi_2}<\pheight{\pi}$. ومن ثم يجري فرض
الاستقراء على~$\pi_1'$ و$\pi_2'$، فتوجد~!!{proof}s~$\pi_1''$ و$\pi_2''$
لـ$\Gamma\Sequent\Delta,!B$ و$\Gamma\Sequent\Delta,!C$، على الترتيب،
مع~$\pheight{\pi_1''}\leq\pheight{\pi_1'}\leq\pheight{\pi_1}$ و
$\pheight{\pi_2''}\leq\pheight{\pi_2'}\leq\pheight{\pi_2}$. وعند ذلك
يكون~!!{proof}~$\pi'$:
\[
\AxiomC{}
\RightLabel{$\pi_1''$}
\Deduce$\Gamma \fCenter \Delta, !B$
\AxiomC{}
\RightLabel{$\pi_2''$}
\Deduce$\Gamma \fCenter \Delta, !C$
\RightLabel{\RightR{\land}}
\BinaryInf$\Gamma \fCenter \Delta, !B \land !C$
\DisplayProof
\]
مع~$\pheight{\pi'}\leq\pheight{\pi}$.

وتحتاج قواعد الكم، حين تكون~!!{formula} رئيسة، إلى تفصيل خاص.

فلتكن~$!A\ident\lforall[x][!B(x)]$ رئيسة في اليمين؛ وعندئذ ينتهي~$\pi$
بـ
\[
\AxiomC{}
\RightLabel{$\pi_1$}
\Deduce$\Gamma \fCenter \Delta, \lforall[x][!B(x)], !B(a)$
\RightLabel{\RightR{\lforall}}
\UnaryInf$\Gamma \fCenter \Delta, \lforall[x][!B(x)], \lforall[x][!B(x)]$
\DisplayProof
\]
وبحسب~\olref{lem:invert-quant} يوجد~!!a{proof}~$\pi_1'$ من~!!{height}
لا يتجاوز~$\pheight{\pi_1}$ لـ$\Gamma\fCenter\Delta,!B(c),!B(a)$ لكل~$c$
لا تظهر في~$\Gamma\fCenter\Delta,\lforall[x][!B(x)],!B(a)$. ونأخذ
$\pi_1'$ منتظمًا؛ فيكون الـ!!{proof}~$\Subst{\pi_1'}{a}{c}$، بموجب
\olref[qua]{lem:subst-regular}، !!a{proof} لـ$\Gamma\fCenter\Delta,!B(a),
!B(a)$، ومن~!!{height} لا يتجاوز~$\pheight{\pi_1}$ كذلك. ثم يمدنا فرض
الاستقراء بـ!!a{proof}~$\pi_1''$ لـ$\Gamma\fCenter\Delta,!B(a)$؛ فإذا
طبقنا~$\RightR{\lforall}$ تم لنا~$\pi'$:
\[
\AxiomC{}
\RightLabel{$\pi_1''$}
\Deduce$\Gamma \fCenter \Delta, !B(a)$
\RightLabel{\RightR{\lforall}}
\UnaryInf$\Gamma \fCenter \Delta, \lforall[x][!B(x)]$
\DisplayProof
\]
مع~$\pheight{\pi'}\leq\pheight{\pi}$.

وأما إن كانت~$!A\ident\lexists[x][!B(x)]$ هي الرئيسة في اليمين، كانت
صورة~$\pi$:
\[
\AxiomC{}
\RightLabel{$\pi_1$}
\Deduce$\Gamma \fCenter \Delta, \lexists[x][!B(x)], \lexists[x][!B(x)], !B(t)$
\RightLabel{\RightR{\lexists}}
\UnaryInf$\Gamma \fCenter \Delta, \lexists[x][!B(x)], \lexists[x][!B(x)]$
\DisplayProof
\]
يجري فرض الاستقراء على~$\pi_1$، فيعطينا~!!a{proof}~$\pi_1'$ لـ$\Gamma
\Sequent\Delta,\lexists[x][!B(x)],!B(t)$، من غير حاجة هنا إلى لمّة عكس.
وعندئذ يكون~!!{proof}~$\pi'$:
\[
\AxiomC{}
\RightLabel{$\pi_1'$}
\Deduce$\Gamma \fCenter \Delta, \lexists[x][!B(x)], !B(t)$
\RightLabel{\RightR{\lexists}}
\UnaryInf$\Gamma \fCenter \Delta, \lexists[x][!B(x)]$
\DisplayProof
\]
مع~$\pheight{\pi'}\leq\pheight{\pi}$.
\end{proof}

\begin{prob}
أنشئ برهان~\olref{prop:G3c-cont-adm} لـ\LeftR{\Contraction}، وفصّل حالتي
$!A\ident(!B\land !C)$ و$!A\ident(!B\lif !C)$.
\end{prob}

\begin{prob}
أنشئ برهان~\olref{prop:G3c-cont-adm} لـ\LeftR{\Contraction} في حالتي الكم
الباقيتين: حين~$!A\ident\lforall[x][!B(x)]$ وحين
$!A\ident\lexists[x][!B(x)]$.
\end{prob}

\end{document}
